\documentclass[11pt,a4paper]{article}
\usepackage[margin=25mm]{geometry}
\usepackage{amsmath,amssymb,amsthm,booktabs}
\setlength{\emergencystretch}{2em}
\newtheorem{theorem}{Theorem}\newtheorem{lemma}[theorem]{Lemma}
\newcommand{\SL}{\operatorname{SL}}\newcommand{\PSL}{\operatorname{PSL}}
\newcommand{\PGL}{\operatorname{PGL}}\newcommand{\PG}{\operatorname{PG}}
\title{Two stabilizers and elementary matrices in dimension three}
\author{Human Compression project}\date{14 October 2026}
\begin{document}\maketitle
We reprove the rank-four construction of Leemans and Vandenschrick
\cite[Theorem 1.1]{source}. Its matrices and duality are retained exactly.
The generation argument below uses elementary matrices; it does not
require a classification of maximal subgroups. We retain the classical
rotation-group correspondence and its rank-four intersection criterion
\cite{SW}, as already stated explicitly in the project's common theory
for projective markings. We also retain the classical description of
automorphisms of $\PSL_3(q)$ as projective semilinear maps and their
compositions with inverse transpose, for $q\geq5$.

Let $F=\mathbb F_q$, where $q\geq5$, and let $x$ generate $F^\times$.
Put
\[
 a=\begin{pmatrix}x&1&0\\0&1+x^{-1}&-x^{-1}\\0&1&0\end{pmatrix},\quad
 b=\begin{pmatrix}-x^{-1}&0&0\\0&0&1\\0&x&0\end{pmatrix},\quad
 c=\begin{pmatrix}x&0&0\\x-x^{-1}&1&0\\x-x^{-1}&0&x^{-1}\end{pmatrix}.
\]
All have determinant one. We use the same letters for their projective
images. Set $d=\gcd(3,q-1)$.
\begin{theorem}
The ordered triple $(a,b,c)$ generates $\PSL_3(q)$ and defines a
self-dual chiral polytope of type
\[
             [q-1,\,2(q-1)/d,\,q-1].
\]
Its two end subgroups are the full set stabilizers of two distinct lines
and two distinct points, respectively. Each has order
$2q^2(q-1)^2/d$.
\end{theorem}

\section{The full stabilizers}
For $u,v\in F$ and $r\in F^\times$, write
\[
 K(u,v)=I+uE_{21}+vE_{31},\qquad h(r)=\operatorname{diag}(r^{-2},r,r).
\]
Thus $K(u,v)K(u',v')=K(u+u',v+v')$ and $b^{2m}=h(x^m)$.
Direct multiplication gives, for $r,s\in F^\times$,
\begin{align}
 h(r)c^{-1}h(s/r)c h(s^{-1})&=K(C/x,C),
       & C&=(r^3-s^3)(1-x^2),\
 c^{-1}K(u,v)c&=K(xu,x^2v). \label{eq:conjugation}
\end{align}
The factor $x^{-1}$ in the first coordinate of the first identity is
essential. The corresponding displayed word on source page 6 omits it.

The $\mathbb F_p$-span $V$ of all $r^3-s^3$ is stable under
multiplication by $x^3$ and contains $x^3-1\ne0$.
Moreover $\mathbb F_p(x^3)=F$: otherwise $x^3$ would belong to a
proper subfield of size at most $\sqrt q$, whereas
\[
 |x^3|=(q-1)/\gcd(3,q-1)\ \geq\ (q-1)/3\ >\ \sqrt q-1.
\]
The strict inequality holds for $q>4$. Hence $V$ is a nonzero vector
space over $F$, so $V=F$. Since $1-x^2\ne0$, the first identity and
addition in $K$ give every $K(t/x,t)$. Conjugation by $c$ gives every
$K(t,x^2t)$. The vectors $(x^{-1},1)$ and $(1,x^2)$ are independent,
with determinant $x-1$. It follows that
\[
                  K(F,F)\ \leq\ \langle b,c\rangle.
\]

Let $p=[0,1,0]^t$, $p'=[0,0,1]^t$. Their full set stabilizer in
$\SL_3(F)$ consists of the matrices
\[
 \begin{pmatrix}\alpha&0&0\\u&\beta&0\\v&0&\gamma\end{pmatrix},
 \quad \alpha\beta\gamma=1,
 \qquad\hbox{and their products with }b.
\]
The subgroup fixing both points is $K(F,F)$ times the determinant-one
diagonal torus. Modulo $K$, the diagonals of $b^{2m}c^n$ have exponent
vectors
\[
                     (-2m+n,\ m,\ m-n).
\]
Every triple of exponents summing to zero occurs: for exponents
$(A,B,C)$ take $m=B$ and $n=A+2B$. Thus $b,c$ generate the full
set stabilizer, even before taking projective images. Dividing its
order by the scalar center gives $2q^2(q-1)^2/d$.

Put
\[
 D=\begin{pmatrix}(x-x^{-1})^{-1}&0&0\\0&1&1\\0&1&x\end{pmatrix},
 \qquad \phi(g)=D(g^{-1})^tD^{-1}.
\]
The matrix $D$ is symmetric and invertible. Consequently $\phi^2=1$,
and direct multiplication gives
\[
             \phi(a)=c^{-1},\qquad
             \phi(b)=b^{-1},\qquad \phi(c)=a^{-1}.
\]
It interchanges $\langle b,c\rangle$ and $\langle a,b\rangle$.
The latter is therefore the full stabilizer of the unordered pair of
lines
\[
                 \ell:y-z=0,\qquad \ell':y-x^{-1}z=0.
\]
One can also identify these lines directly from the actions of $a,b$;
$a$ fixes them and $b$ exchanges them.

\section{Generation and intersection}
Let $\widetilde G=\langle a,b,c\rangle\leq\SL_3(F)$.
The preceding argument gives all $I+tE_{21}$ and $I+tE_{31}$ in
$\widetilde G$. Since $\phi$ preserves $\widetilde G$, it also gives
the images of all these matrices under $\phi$. The block diagonal
shape of $D$ shows that these images range over all
\[
                         I+uE_{12}+vE_{13}.
\]
Indeed inverse transpose changes $K(u,v)$ into
$I-uE_{12}-vE_{13}$, and conjugation by $D$ applies an invertible
linear map to $(u,v)$. Commutators now give the remaining root groups:
\[
 [I+sE_{21},I+tE_{13}]=I+stE_{23},\qquad
 [I+sE_{31},I+tE_{12}]=I+stE_{32},
\]
where $[g,h]=ghg^{-1}h^{-1}$. These six elementary root groups generate
$\SL_3(F)$ by row elimination. This proves generation without any
overgroup classification.

Write $H=\langle b,c\rangle$ and $A=\langle a,b\rangle$ in the
projective group. An element of $A\cap H$ preserves both unordered
pairs $\{p,p'\}$ and $\{\ell,\ell'\}$. It fixes
$e=[1,0,0]^t=\ell\cap\ell'$, and hence has a diagonal or a
monomial representative
\[
 \operatorname{diag}(\alpha,\beta,\gamma),\qquad
 \begin{pmatrix}\alpha&0&0\\0&0&\beta\\0&\gamma&0\end{pmatrix}.
\]
On the line $pp'$, the two distinguished intersections with $\ell,\ell'$
have ratios $z/y=1,x$. A diagonal representative preserves this pair
only when $\gamma/\beta=1$, since $x^2\ne1$. The determinant then
makes it $h(\beta)$, a power of $b^2$. A monomial representative
sends $z/y$ to $\gamma/(\beta(z/y))$; pair preservation forces
$\gamma/\beta=x$. Its determinant makes it $h(\beta)b$, again a
power of $b$. Thus $A\cap H=\langle b\rangle$.

Every power of $c$ fixes $p,p'$ separately. An element shared with
$\langle b\rangle$ must therefore be an even power of $b$, acting
as a scalar on $pp'$. But $c^n$ acts there with eigenvalue ratio
$x^n$, so it is shared only if $q-1$ divides $n$. The eigenvalues of
$c$ are the distinct elements $x,1,x^{-1}$; hence its projective
order is $q-1$ and this shared element is the identity. Applying
$\phi$ gives the other adjacent cyclic intersection. We have proved
\[
 \langle a\rangle\cap\langle b\rangle=1,
 \quad \langle b\rangle\cap\langle c\rangle=1,
 \quad A\cap H=\langle b\rangle.
\]
Direct multiplication gives $(ab)^2=(bc)^2=(abc)^2=I$;
each product is nonscalar. The retained rank-four criterion now proves
the full rotation intersection property.

\section{Excluding reversal}
Suppose an automorphism $\alpha$ of $\PSL_3(F)$ sends $a$ to $a^{-1}$
and fixes $ab$ and $abc$. Then it fixes $c=(ab)^{-1}(abc)$ and
preserves $A=\langle a,ab\rangle$.

The group $H$ fixes no projective point. Indeed its subgroup $K(F,F)$
fixes precisely the points on $pp'$; the full diagonal torus has only
$p,p'$ as common fixed points there, and $b$ exchanges them. Dually,
$A$ fixes no line, although it fixes $e$. A graph automorphism followed
by a semilinear collineation would take $A$ to a group fixing a line.
It therefore cannot preserve $A$. By the retained automorphism theorem,
$\alpha$ is induced by a semilinear collineation with field part
$t\mapsto t^{p^i}$, $0\leq i<[F:\mathbb F_p]$.

The axes of $ab$ and $abc$ are respectively
\[
                        z=0,\qquad (1-x)y_1=z,
\]
where $y_1$ denotes the first homogeneous coordinate.
These follow, in both characteristics, from their displayed matrices;
their axes are uniquely determined by the involutions. Their intersection
is $p$, so the collineation fixes $p$. It also permutes the three
eigenpoints of $a$, namely
\[
 [1,0,0]^t,\quad[1,1-x,1-x]^t,\quad
 [1,x^{-1}-x,1-x^2]^t.
\]
Only the first lies on the invariant line $z=0$, so it fixes $e$.
The linear part $P$ consequently has the form
\[
                     P=\begin{pmatrix}u&0&w\\0&v&t\\0&0&s\end{pmatrix}.
\]
Projective centralization of $c$ means
$Pc^{(p^i)}P^{-1}=\mu c$ for a scalar $\mu$.
Applying this equation to $p$, whose eigenvalue is one on both sides,
gives $\mu=1$. The other two eigenvalues then imply
$x^{p^i}=x$ or $x^{-1}$. The second option is impossible:
for $0\leq i<[F:\mathbb F_p]$ and $q\geq5$,
$p^i+1<q-1$. The first option implies $i=0$, since $x$ is primitive.
The three eigenpoints of $c$ must therefore be individually fixed.
They are $p,p'$ and $[x,1+x,x]^t$. Fixing $p'$ gives $w=t=0$;
fixing the last point gives $u=v=s$, as $x(1+x)\ne0$.
Thus $P$ is scalar and $\alpha$ is the identity. This contradicts
$\alpha(a)=a^{-1}$, because $a$ has order $q-1>2$.
The marking is therefore chiral by the retained correspondence.

Finally, $a,c$ both have order $q-1$. Odd powers of $b$ interchange
$p,p'$ and cannot be scalar; $b^{2m}=h(x^m)$ is scalar precisely
when $x^{3m}=1$. Hence $|b|=2(q-1)/d$. The automorphism $\phi$
sends the marking to its dual $(c^{-1},b^{-1},a^{-1})$, proving
self-duality and the claimed type.

\section{Attribution and scope}
The matrices, stabilizer descriptions, duality and rank-four existence
result are prior work of the source authors. We supply the finite-field
span argument, the full elementary generation proof and the detailed
projective scalar check in the reversal argument. The source's displayed
unipotent word needs the factor $x^{-1}$ noted above. Its next-page
transitivity paragraph also interchanges the two named end subgroups;
the elementary generation argument avoids that paragraph entirely.
These corrections do not change Theorem 1.1.

This component treats Theorem 1.1 and its proof, not the source's rank-five
construction, subgroup classifications or global rank obstruction.
It gives no new answer to an open question. The automorphism theorem
and the polytope correspondence remain explicit classical premises;
finite examples are controls on the formulas, not the proof of the
all-field statement.
\begin{thebibliography}{3}
\bibitem{source} D. Leemans and A. Vandenschrick,
\emph{On chiral polytopes having a group $\PSL(3,q)$ as automorphism
group}, fixed 25-page author edition, 2022, Theorem 1.1 and Section 3.
\bibitem{SW} E. Schulte and A. I. Weiss, \emph{Chiral polytopes},
Applied Geometry and Discrete Mathematics, DIMACS 4 (1991), 493--516;
rotation correspondence and rank-four intersection criterion.
\bibitem{common} Human Compression project,
\emph{Common theory for projective markings and rotations},
Topic E, version 3, Section 2; retained rotation correspondence.
\end{thebibliography}
\end{document}
